%FIXME: Figure out exact placement of cover image

\pagestyle{empty}

% Cover image
\changepage{4.9cm}{4cm}{}{-2cm}{}{}{}{}{}
\definecolor{fenicsred}{RGB}{215, 30, 40}
\pagecolor{fenicsred}
\hfill
\includegraphics{frontmatter/eps/editors.eps} \\
\includegraphics[height=8cm]{frontmatter/eps/title.eps}
\vfill
\hfill
\includegraphics[height=12cm]{frontmatter/eps/f.eps}
\thispagestyle{empty}
\newpage
\changepage{-4.9cm}{-4cm}{}{2cm}{}{}{}{}{}
\pagecolor{white}

% Some useful information
\null\vfill
\noindent
\begin{minipage}{\textwidth}
  Automated Scientific Computing \\
  Logg, Mardal, Wells (editors) \\ \\
  This is a preliminary draft, \today. \\
  Send comments and suggestions to \texttt{fenics-book-dev@fenics.org}.
\end{minipage}
\newpage

% Title page
\begin{center}

  \vspace{1cm}

  \LARGE

  Automated Scientific Computing

  \vspace{1cm}

  \emph{The FEniCS Book}

  \vfill

  \Large

  Preliminary draft, \today

\end{center}
\cleardoublepage

% Nice qute
\null\vspace{5cm}
\linespread{1.5}
\begin{center}
  \begin{quote}
    \it

    And now that we may give final praise to the machine we may say that
    it will be desirable to all who are engaged in computations which, it
    is well known, are the managers of financial affairs, the
    administrators of others' estates, merchants, surveyors, geographers,
    navigators, astronomers\ldots For it is unworthy of excellent men to lose
    hours like slaves in the labor of calculations which could safely be
    relegated to anyone else if the machine were used.
    
    \rm

    \vspace{0.5cm}
    
    Gottfried Wilhelm Leibniz (1646--1716)
  \end{quote}
\end{center}
\linespread{1.0}
\cleardoublepage

\pagestyle{fancy}
